\section{Introduction}
\label{sec:introduction}

\IEEEPARstart{S}{pectrum} scarcity has become a defining bottleneck for
modern and forthcoming wireless systems. Dynamic spectrum access (DSA),
in which secondary users (SUs) opportunistically reuse spectrum left idle by
primary users (PUs), is widely regarded as a key remedy and underpins
emerging shared-spectrum regimes such as the citizen broadband radio service
(CBRS)~\cite{sohul2015cbrs}. Realizing DSA requires solving two tightly
coupled subproblems: \emph{spectrum sensing}, which estimates what spectrum
is free, and \emph{spectrum access}, which decides how to use it. Because the
quality of sensing directly bounds what access can achieve, and because the
performance after access is the only meaningful measure of whether a sensing
decision was good, a joint treatment of the two is necessary rather than
optional.

Two largely separate research lines have matured around these subproblems.

The first, \emph{time--frequency (TF) localization} (TFL), treats the
wideband spectrogram as an image and applies object detection to output, for
every emission, a bounding box describing its center frequency, bandwidth,
start time, and duration~\cite{prasad2020,li2022hf}. Because the wideband
spectrogram and image object detection are structurally similar, a sequence
of detectors has been adapted to TFL, including region-based and You Only Look
Once (YOLO)-family architectures~\cite{prasad2020,li2022hf}, anchor-free
keypoint detectors that resolve the prior-anchor mismatch caused by the diverse
aspect ratios of radio-frequency (RF) emissions~\cite{zhao2024uav}, and transformer backbones that remain robust
when emissions overlap in the TF plane~\cite{zhao2025trtfl_ssl,
zhang2024spectrumtransformer}. This line has steadily improved how precisely
and robustly a wideband sensor can describe occupancy. Yet it terminates at
sensing: it produces a rich, continuous, often uncertainty-bearing
description of the spectrum and then does nothing with it. The access
decision is left out of scope.

The second line applies \emph{deep reinforcement learning (DRL) to DSA},
learning access---and sometimes sensing---policies under partial
observability~\cite{wang2018,zhong2019,xu2020partial,li2020aggregation,
li2024hmadrl,bokobza2023ddqsa}. Deep $Q$-networks (DQNs)~\cite{wang2018},
actor--critic methods~\cite{zhong2019}, recurrent $Q$-networks for partial
observation~\cite{xu2020partial,hausknecht2015drqn}, and channel-aggregation
schemes~\cite{li2020aggregation} have progressively enlarged the action space
and improved sample efficiency. Almost without exception, however, this line
assumes the wideband spectrum is pre-partitioned into a \emph{fixed grid of
channels}, each summarized by a binary good/bad or idle/busy state.

These two lines are separated by an assumption that is rarely examined. The
DRL-DSA line inherits a fixed channelization in which PU activity is aligned
to grid cells. But real wideband emissions occupy \emph{arbitrary} center
frequencies, bandwidths, and durations---and in shared-spectrum regimes such
as CBRS the center frequency and subchannel division of incumbents change
adaptively with traffic and channel state~\cite{sohul2015cbrs,zhao2025trtfl_ssl}.
This is precisely the regime the TFL line was built to handle. Once occupancy
is not grid-aligned, the ``good/bad channel'' state that DRL-DSA agents
consume is itself a lossy, and sometimes incorrect, summary of the
environment: a single grid cell may straddle the boundary of a PU emission,
forcing the agent either to collide or to abandon usable spectrum. To our
knowledge, no prior work feeds the \emph{continuous, uncertainty-bearing} TF
occupancy estimate produced by a localizer \emph{into} the access decision as
the agent's state, nor analyzes how the localizer's errors propagate to access
performance and PU protection. Closing that loop, and analyzing it, is the
subject of this paper.

We propose \emph{Localization-Aware Joint Spectrum Sensing and Access}
(LA-JSSA), in which the sensing module emits a continuous occupancy
\emph{belief map} over the TF plane that directly serves as the access
agent's state, and the agent selects a continuous TF resource block for
transmission. We further develop a distributional~\cite{dabney2018qrdqn}, risk-sensitive access tier
(D-LA-JSSA) tailored to the belief-state representation. Our contributions
are:
\begin{enumerate}
\item \textbf{A belief-map closed loop.} We formulate joint sensing and
access as a partially observed control problem whose state is the continuous TF occupancy
posterior rather than a fixed-channel vector
(Section~\ref{sec:system_model}), realized by a two-timescale architecture
that couples a supervised TFL sensing tier to a reinforcement-learned access
tier without unstable end-to-end training (Section~\ref{sec:framework}).
\item \textbf{Error-propagation analysis and a three-way trade-off.} We map
cell-level localization error ($\Pmd,\Pfa$) to PU-collision probability and
SU throughput, yielding a sensing-accuracy/throughput/PU-protection
trade-off with a characterized optimal belief threshold $\thr^\star$
(Theorem~\ref{thm:tradeoff}), and a linear bound on the price of imperfect
localization (Corollary~\ref{cor:price}).
\item \textbf{Risk-sensitive access via a distributional critic.} We cast
access as a \emph{distributional} reinforcement-learning (RL)
problem~\cite{dabney2018qrdqn,ma2020dsac} with a tail-risk conditional
value-at-risk (CVaR) constraint on PU collision, and show its operating point is at least as
conservative as the mean-constrained one, with a first-order estimate of the
conservativeness margin (Theorem~\ref{thm:cvar},
Corollary~\ref{cor:cvargap}).
\item \textbf{Information refinement and a stylized regret analysis.} We show
the belief-state policy class is no worse than any coarser fixed-channel
representation under a common action space
(Proposition~\ref{thm:sufficiency}), with a strict gap on a misaligned
instance (Example~\ref{ex:strictgap}); and, for a tractable tabular
surrogate, we separate the regret into a vanishing learning term and a
policy-independent sensing gap (Proposition~\ref{thm:regret}). The latter is a
stylized analysis, not a guarantee for the deep agent.
\item \textbf{Validation.} Simulations on a continuous-bandwidth, non-grid
PU model confirm the predicted operating point and the tail-risk reduction,
validate the full sensing-error propagation chain end to end, and show the
belief-map advantage is monotone in representation richness and robust across
traffic, geometry, and SNR regimes (Section~\ref{sec:simulation}).
\end{enumerate}
The remainder of the paper is organized as follows.
Section~\ref{sec:related_work} reviews related work,
Section~\ref{sec:system_model} formulates the model,
Section~\ref{sec:framework} presents LA-JSSA,
Section~\ref{sec:theory} develops the theory,
Section~\ref{sec:simulation} reports experiments, and
Section~\ref{sec:conclusion} concludes.
